\chapter{Introduction}
\label{ch:introduction}

\section{Motivation}

Robotic intercropping is an emerging agricultural strategy that uses autonomous robots to cultivate multiple crop species in the same field. Field conditions constantly change, and treating the environment as constant can result in poor navigation, unnecessary stops, or damage to plants\cite{RoboticIntercroppingProject2026}. Therefore, this project aims to equip the robot with the ability to recognise the type of terrain that it experiences while driving, using only the available intrinsic motion sensors. Specifically, accelerometer and gyroscope measurements, together with wheel odometry, are used to capture vibration patterns, impacts, and motion signatures arising from robot–terrain interaction, including transitions between surface types.

The study considers two complementary terrain scenarios. Before the main data collection, a campus route covering smooth pavement, grass, cobblestone, and a dirt track (recorded under both dry and wet/muddy conditions on different days) served as the development environment to build and test the full processing pipeline. The campus route was chosen to mimic different terrains found on the farm, while also including classes that provide a deliberately wider roughness contrast, notably smooth asphalt and cobblestone, which bracket the farm classes from above (hard, impulsive) and below (smooth, low-vibration). The clear distinction between asphalt/cobblestone and the softer surfaces also provides a useful reference point for comparison with the primary scenario: an actual farm, where grass, soil, and rough dirt tracks present a harder classification problem due to their greater similarity. Data were collected across multiple repeated runs on each route, four campus runs on separate days and five farm runs across two days, to assess cross-run generalisation. The effect of driving speed was also explicitly investigated by evaluating models on both the full recorded speed range and a lower, more consistent speed regime, with the latter strategy adopted during the main farm data collection. The collected data were labelled by terrain segment, and the project developed and evaluated methods ranging from classical machine-learning baselines to a 1D convolutional neural network. Transition detection between adjacent segments was examined as a harder sub-problem.

Overall, the project does not rely on external perception sensors such as cameras or LiDAR, but instead investigates the limits of proprioceptive terrain recognition. The goal is not to replace exteroceptive sensing, but to explore how intrinsic motion signals can complement it, providing a lightweight, always-available channel for terrain awareness that does not depend on lighting conditions, sensor range, or additional hardware.

\newpage
\section{BSc Learning Objectives}

This project addresses the five DTU BSc graduate competencies \cite{DTU_BachelorProjectRules2026} as summarised in \Cref{tab:learning_objectives}.

\begin{table}[H]
\centering
\small
\renewcommand{\arraystretch}{1.25}
\begin{tabular}{p{0.30\textwidth}p{0.62\textwidth}}
\toprule
\textbf{DTU BSc Objective} & \textbf{How it is met in this project} \\
\midrule
1. Structure and execute a major project & Planned and executed an 18-week project with weekly supervision, defined milestones, and timely intermediate deliverables (frozen dataset, baseline results, draft report). \\

2. Summarise and interpret technical information & Synthesised literature on terrain classification, IMU-based event detection, time-series feature extraction, and cross-run model evaluation; interpreted results in relation to robotic-intercropping requirements. \\

3. Work across all project phases & Formulated the problem statement, designed the experimental protocol, implemented the data-processing pipeline and ML models, and produced reproducible technical documentation. \\

4. Independently acquire new knowledge and apply it critically & Independently studied and applied time- and frequency-domain features, mutual-information feature selection, classical and deep classifier families, and Leave-One-Run-Out cross-validation as they were needed during implementation; and iteratively refined the project itself: a scope pivot from collision detection to terrain classification, and protocol adjustments informed by the campus campaign. \\

5. Communicate in written, visual, and oral form & Reported findings in the present report, in figures (confusion matrices, MI rankings, sensor-ablation plots), and in oral status updates and the final presentation. \\
\bottomrule
\end{tabular}
\caption{Mapping of the project's activities to the DTU BSc graduate competencies.}
\label{tab:learning_objectives}
\end{table}

\section{Problem Statement}
This thesis addresses the following question:
\begin{quote}
\emph{Can terrain type be reliably classified from accelerometer, gyroscope, and wheel-odometry data collected by a wheeled agricultural robot, without any external perception sensor?}
\end{quote}
The question decomposes into four sub-problems:
\begin{enumerate}
    \item Designing and executing a repeatable data-collection and preprocessing pipeline that turns raw sensor logs from both a controlled campus development route and an operational farm into a labelled, model-ready feature table.
    \item Training and comparing multiple classifier families, from classical baselines to a 1D convolutional neural network, under a Leave-One-Run-Out cross-validation scheme that prevents leakage between training and test runs.
    \item Characterising how driving speed and terrain similarity affect classification, by evaluating each model on both the full and a low-constant speed regime, and by contrasting the campus and farm scenarios under the same pipeline.
    \item Examining transition detection between adjacent terrain segments as a harder sub-problem than steady-state classification.
\end{enumerate}


\section{Project Scope and Delimitations}

The scope of this thesis is bounded as follows. \emph{Sensors:} only the on-board IMU (accelerometer, gyroscope) and wheel-encoder odometry are used; cameras, LiDAR, GNSS, the magnetometer, and the gyroscope's yaw channel (which encodes steering rather than terrain) are excluded. \emph{Platform:} all data are collected with a single robot, \emph{Ricbot}, in its differential-drive configuration; generalisation to other chassis is not claimed. \emph{Terrain classes:} only the surfaces encountered in the two campaigns are considered: campus (asphalt, grass, cobblestone, dry and muddy dirt track) and farm (grass, tilled soil, compacted dirt track); rare events such as deep puddles and steep slopes are out of scope. \emph{Operating conditions:} the farm dataset was recorded entirely under dry, sunny weather and at low, near-constant speed; wet-soil operation and high-speed regimes are not evaluated. \emph{Cross-domain transfer:} the campus and farm datasets are treated as two independent classification problems with their own label spaces, because the labels only partially overlap. \emph{Deployment:} classification is performed offline on pre-recorded logs; real-time and embedded deployment are left as future work.


\section{Project Plan and Execution}
\label{sec:project-plan}

The project was scoped as an 18-week effort (2 February – 8 June 2026), organised into six phases from initial planning through to final submission. The original weekly plan and milestones are summarised below; the resulting Gantt chart is shown in \Cref{fig:gantt}.

\subsection{Planned Schedule}

\begin{itemize}
    \item \textbf{Phase 0 — Project Management (W1–2):} define research questions, terrain classes, and success criteria; risk assessment and timeline.
    \item \textbf{Phase 1 — Data Foundation (W2–3):} ingestion pipeline with time alignment, automated quality control, and a labelling system.
    \item \textbf{Phase 2 — Proxy Data \& Preprocessing (W2–7):} controlled proxy data collection, resampling and sliding-window extraction, and the first time- and frequency-domain feature set; \emph{Proxy Dataset Freeze}.
    \item \textbf{Phase 3 — Modelling, Evaluation \& Field Collection (W8–11):} classical baselines (LR, RF, XGBoost, SVM) and neural models (MLP, 1D-CNN) under Leave-One-Run-Out cross-validation, followed by outdoor field data collection and a proxy-to-field domain transfer evaluation.
    \item \textbf{Phase 4 — Validation \& Documentation (W12–15):} robustness tests (noise injection, sensor dropout) and the methodology and results writing sprint.
    \item \textbf{Phase 5 — Finalisation (W16–18):} introduction, literature review, revisions, and submission on 8 June 2026.
\end{itemize}


\definecolor{phase0}{RGB}{80, 80, 80}
\definecolor{phase1}{RGB}{120, 120, 120}
\definecolor{phase2}{RGB}{0, 102, 204}
\definecolor{phase3}{RGB}{0, 153, 153}
\definecolor{phase4}{RGB}{204, 102, 0}
\definecolor{phase5}{RGB}{153, 0, 0}
\begin{figure}[htbp]
    \centering
    \makebox[\textwidth][c]{%
    \resizebox{1\textwidth}{!}{%
    \begin{ganttchart}[
        canvas/.style={draw=black!25, line width=.6pt},
        hgrid style/.style={draw=black!25, line width=.6pt},
        vgrid={*1{draw=black!5, line width=.6pt}},
        title/.style={draw=black!25, fill=none},
        title label font=\bfseries\normalsize,
        include title in canvas=false,
        bar/.append style={draw=none},
        bar height=.6,
        y unit title=1.2cm,
        y unit chart=1.1cm,
        bar label font=\normalsize\color{black!80},
        bar label node/.append style={align=right, text width=6cm},
        group label font=\normalsize\bfseries\color{black!85},
        group label node/.append style={align=right, text width=6cm},
        progress label text={},
        group right shift=0,
        group top shift=.6,
        group height=.3,
        x unit=0.34cm,
        time slot format=isodate,
        calendar week text={W\currentweek},
        link/.style={-stealth, draw=black!45, rounded corners=2pt}
    ]{2026-02-02}{2026-06-12}

        \gantttitlecalendar{month=name, week=1} \\

        % PHASE 0
        \ganttgroup[group/.append style={fill=phase0}]{Phase 0: Project Management}{2026-02-02}{2026-02-14} \\
        \ganttbar[bar/.append style={fill=phase0}, progress=100, name=P0A]{Research Questions \& Metrics}{2026-02-02}{2026-02-07} \\
        \ganttbar[bar/.append style={fill=phase0}, progress=100, name=P0B]{Risk Assessment \& Project Plan}{2026-02-08}{2026-02-14} \\

        % PHASE 1
        \ganttgroup[group/.append style={fill=phase1}]{Phase 1: Data Foundation}{2026-02-15}{2026-02-28} \\
        \ganttbar[bar/.append style={fill=phase1}, progress=100, name=P1A]{Pipeline \& QC Framework}{2026-02-15}{2026-02-21} \\
        \ganttbar[bar/.append style={fill=phase1}, progress=100, name=P1B]{Labelling System (Provisional)}{2026-02-22}{2026-02-28} \\

        % PHASE 2
        \ganttgroup[group/.append style={fill=phase2}]{Phase 2: Proxy Data \& Preprocessing}{2026-02-15}{2026-03-21} \\
        \ganttbar[bar/.append style={fill=phase2}, name=P2A]{Proxy Data Collection}{2026-02-15}{2026-03-07} \\
        \ganttbar[bar/.append style={fill=phase2}, name=P2C]{Sync, Resample \& Windowing}{2026-03-08}{2026-03-14} \\
        \ganttbar[bar/.append style={fill=phase2}, name=P2D]{Feature Engineering Foundation}{2026-03-15}{2026-03-21} \\
        \ganttmilestone[milestone/.append style={fill=phase2}, name=MS1]{Proxy Dataset Freeze}{2026-03-21} \\

        % PHASE 3
        \ganttgroup[group/.append style={fill=phase3}]{Phase 3: Modelling \& Evaluation}{2026-03-22}{2026-04-25} \\
        \ganttbar[bar/.append style={fill=phase3}, name=P3A]{Baseline Models (RF/XGB/LR)}{2026-03-22}{2026-03-28} \\
        \ganttbar[bar/.append style={fill=phase3}, name=P3B]{Neural Network (MLP/1D-CNN)}{2026-03-29}{2026-04-04} \\
        \ganttbar[bar/.append style={fill=phase3}, name=P3C]{Field Data Collection (Outdoor)}{2026-04-05}{2026-04-11} \\
        \ganttbar[bar/.append style={fill=phase3}]{Field Data Collection (If more data is required -  Optional)}{2026-04-19}{2026-04-25} \\
        \ganttbar[bar/.append style={fill=phase3}, name=P3D]{Domain Transfer \& Evaluation}{2026-04-12}{2026-04-18} \\
        \ganttmilestone[milestone/.append style={fill=phase3}, name=MS2]{Final Model Selection}{2026-04-18} \\

        % PHASE 4
        \ganttgroup[group/.append style={fill=phase4}]{Phase 4: Validation \& Documentation}{2026-04-19}{2026-05-16} \\
        \ganttbar[bar/.append style={fill=phase4}, name=P4A]{Robustness \& Noise Tests}{2026-04-19}{2026-05-02} \\
        \ganttbar[bar/.append style={fill=phase4}, name=P4B]{Methodology \& Results Sprint}{2026-05-03}{2026-05-16} \\

        % PHASE 5
        \ganttgroup[group/.append style={fill=phase5}]{Phase 5: Finalisation}{2026-05-17}{2026-06-08} \\
        \ganttbar[bar/.append style={fill=phase5}, name=P5A]{Lit Review \& Conclusions}{2026-05-17}{2026-05-30} \\
        \ganttbar[bar/.append style={fill=phase5}, name=P5B]{Final Revisions \& Submission}{2026-05-31}{2026-06-07} \\
        \ganttmilestone[milestone/.append style={fill=phase5}, name=MS3]{FINAL SUBMISSION}{2026-06-08}

        % Links
        \ganttlink{P0A}{P0B}
        \ganttlink{P0B}{P1A}
        \ganttlink{P1A}{P1B}
        \ganttlink{P2A}{P2C}
        \ganttlink{P2C}{P2D}
        \ganttlink{P2D}{P3A}
        \ganttlink{P3A}{P3B}
        \ganttlink{P3B}{P3C}
        \ganttlink{P3C}{P3D}
        \ganttlink{P3D}{P4A}
        \ganttlink{P4A}{P4B}
        \ganttlink{P4B}{P5A}
        \ganttlink{P5A}{P5B}

    \end{ganttchart}
    }}
    \caption{Originally planned project timeline (2 February – 8 June 2026).}
    \label{fig:gantt}
\end{figure}
\vspace{-4mm}
\begin{figure}[H]
    \centering
    \makebox[\textwidth][c]{%
    \resizebox{1\textwidth}{!}{%
    \begin{ganttchart}[
        canvas/.style={draw=black!25, line width=.6pt},
        hgrid style/.style={draw=black!25, line width=.6pt},
        vgrid={*1{draw=black!5, line width=.6pt}},
        title/.style={draw=black!25, fill=none},
        title label font=\bfseries\normalsize,
        include title in canvas=false,
        bar/.append style={draw=none},
        bar height=.6,
        y unit title=1.2cm,
        y unit chart=1.1cm,
        bar label font=\normalsize\color{black!80},
        bar label node/.append style={align=right, text width=6cm},
        group label font=\normalsize\bfseries\color{black!85},
        group label node/.append style={align=right, text width=6cm},
        progress label text={},
        group right shift=0,
        group top shift=.6,
        group height=.3,
        x unit=0.34cm,
        time slot format=isodate,
        calendar week text={W\currentweek},
        link/.style={-stealth, draw=black!45, rounded corners=2pt}
    ]{2026-02-02}{2026-06-12}

        \gantttitlecalendar{month=name, week=1} \\

        % PHASE 0 — as planned
        \ganttgroup[group/.append style={fill=phase0}]{Phase 0: Project Management}{2026-02-02}{2026-02-14} \\
        \ganttbar[bar/.append style={fill=phase0}, progress=100]{Research Questions \& Scope Pivot}{2026-02-02}{2026-02-14} \\

        % PHASE 1 — slipped ~4 days
        \ganttgroup[group/.append style={fill=phase1}]{Phase 1: Data Foundation}{2026-02-15}{2026-03-04} \\
        \ganttbar[bar/.append style={fill=phase1}, progress=100]{Pipeline, QC \& Labelling}{2026-02-15}{2026-03-04} \\

        % PHASE 2 — campus data ran longer
        \ganttgroup[group/.append style={fill=phase2}]{Phase 2: Campus Data \& Pipeline}{2026-02-23}{2026-03-27} \\
        \ganttbar[bar/.append style={fill=phase2}, progress=100]{Campus Data Collection (4 runs)}{2026-02-23}{2026-03-26} \\
        \ganttbar[bar/.append style={fill=phase2}, progress=100]{Sync, Resample \& Features}{2026-03-05}{2026-03-19} \\
        \ganttmilestone[milestone/.append style={fill=phase2}]{Campus Dataset Freeze}{2026-03-27} \\

        % PHASE 3 — extended; absorbs Phase 4
        \ganttgroup[group/.append style={fill=phase3}]{Phase 3: Modelling \& Analysis}{2026-03-22}{2026-05-11} \\
        \ganttbar[bar/.append style={fill=phase3}, progress=100]{Classical + Neural Models (Campus)}{2026-03-22}{2026-04-10} \\
        \ganttmilestone[milestone/.append style={fill=phase3}]{Speed-Regime A/B Split Added}{2026-03-30} \\
        \ganttbar[bar/.append style={fill=phase3}, progress=100]{Sensor Ablation \& Transition (Campus)}{2026-04-11}{2026-04-20} \\
        \ganttbar[bar/.append style={fill=phase3}, progress=100]{Farm Data Collection (5 runs)}{2026-04-21}{2026-04-28} \\
        \ganttbar[bar/.append style={fill=phase3}, progress=100]{Farm Models + 3-class Matched Analysis}{2026-04-29}{2026-05-05} \\
        \ganttbar[bar/.append style={fill=phase3}, progress=100]{CUSUM Change Detection}{2026-05-06}{2026-05-09} \\
        \ganttmilestone[milestone/.append style={fill=phase3}]{Final Model Comparison Locked}{2026-05-11} \\

        % PHASE 5 — writing; started ~2 weeks earlier than planned
        \ganttgroup[group/.append style={fill=phase5}]{Phase 5: Writing \& Submission}{2026-05-04}{2026-06-08} \\
        \ganttbar[bar/.append style={fill=phase5}, progress=100]{Methodology \& Results}{2026-05-04}{2026-05-16} \\
        \ganttbar[bar/.append style={fill=phase5}, progress=100]{Introduction \& Lit Review}{2026-05-17}{2026-05-30} \\
        \ganttbar[bar/.append style={fill=phase5}]{Final Revisions \& Submission}{2026-05-31}{2026-06-07} \\
        \ganttmilestone[milestone/.append style={fill=phase5}]{FINAL SUBMISSION}{2026-06-08}

    \end{ganttchart}
    }}
    \caption{Actual project execution timeline. Phase~4 (Validation \&
    Documentation) was absorbed into Phase~3 as the sensor-ablation
    and transition-detection studies. The writing phase started on
    4 May, roughly two weeks ahead of the original plan, to absorb
    the slip in the final model-comparison milestone.}
    \label{fig:gantt_actual}
\end{figure}


\subsection{Deviations from the Plan}

The overall six-phase structure held, but several elements shifted
(\Cref{fig:gantt_actual}):

\begin{itemize}
    \setlength{\itemsep}{2pt}
    \item \textbf{Scope pivot during planning.} The original target was terrain-induced collision detection; this was dropped because collecting sufficient collision data was  infeasible without risking damage to Ricbot, and the scope was refocused on terrain-segment classification.
    \item \textbf{Data collection ran longer than planned} (campus: 23 Feb – 26 Mar instead of 15 Feb – 7 Mar; farm: 21 \& 28 Apr instead of 5–11 Apr), driven by robot availability, maintenance and weather. The campus dataset freeze slipped from 21 March to 27 March.
    \item \textbf{Speed-regime A/B split added (30 March).} Early campus results showed that speed variability and class imbalance were major confounds, so an explicit comparison between the full speed range (Dataset~A) and a low, near-constant regime (Dataset~B) was added. The low-speed strategy was subsequently adopted on the farm, a more realistic deployment scenario.
    \item \textbf{Campus dataset role shifted; domain transfer evaluation dropped.} The campus campaign was originally intended both to develop the pipeline and to feed a direct train-on-campus / test-on-farm transfer evaluation. The farm classes turned out to have no clean one-to-one campus analogues, so the campus surfaces were instead selected as a deliberately wider roughness range: classes plausibly found on a farm (grass, dirt track) together with deliberately extreme contrasts (smooth asphalt, cobblestone) chosen to probe how the pipeline behaves at the easy and hard ends of the recognition problem. The two campaigns are therefore reported as independent classification problems with their own label spaces.
     \item \textbf{Additional analyses added (18 April – 9 May):} sensor-zeroing and the implicit $n_{\text{stable}}$ transition detector as new evaluation axes; a three-class matched campus subset for a like-for-like comparison with the farm; and, late in Phase~3, an explicit CUSUM change detector on the classifier posteriors, added on the supervisor's recommendation as a principled alternative with a calibratable false-alarm rate, following \cite[\S6.4]{Blanke2003}.
    \item \textbf{Phase~4 absorbed into Phase~3; writing started two weeks earlier.} The ablation work replaced the originally separate Phase~4. The final model comparison was locked on 11 May (vs.\ planned 18 April) due to the farm data collection delay and some pipeline changes, and writing began on 4 May (vs.\ planned 17 May).
\end{itemize}


\section{Thesis Outline}

\Cref{ch:background} reviews proprioceptive terrain classification, IMU signal physics, time- and frequency-domain features, supervised learning, the velocity--vibration confound and change-point detection. \Cref{ch:setup,ch:collection} describe the Ricbot platform and sensor suite, the two outdoor campaigns (campus development and primary farm) and the video-aided labelling procedure. The methodology spans \Cref{ch:processing} (QC, synchronisation, cleaning, filtering and sliding-window segmentation), \Cref{ch:feature} (189 time- and frequency-domain features, the speed-regime A/B split and MI ranking with correlation pruning) and \Cref{ch:classification} (six classifier families under Leave-One-Run-Out cross-validation). \Cref{ch:results} reports classification, sensor-ablation and transition-detection results on the farm and three nested campus configurations; \Cref{ch:discussion,ch:conclusion} interpret them and outline future work.

In brief, intrinsic motion signals alone suffice to recognise terrain in both settings, careful feature engineering keeps classical learners competitive with a 1D-CNN, and the accelerometer and gyroscope carry the bulk of the discriminative information.